% Part: proof-theory
% Chapter: cut-elimination
% Section: midsequent

\documentclass[../../../include/open-logic-section]{subfiles}

\begin{document}

\olfileid{pt}{cut}{mid}

\olsection{في مبرهنة المتتالية الوسطى ومبرهنة هيربراند}

ومن أجلّ ما يترتب على حذف القطع مبرهنة المتتالية الوسطى. ومؤداها أنه إذا
وجد~!!a{proof}~$\pi$ لـ${\OLId{greek-variable}{Gamma}} \Sequent {\OLId{greek-variable}{Delta}}$، وكانت كل !!{formula}
فيه سوَرية المقدمة، وجد~!!a{proof}~$\pi'$ يشتمل على متتالية
${\OLId{latin-looped}{S}} \ident \Pi \Sequent \Lambda$ تسمى \emph{المتتالية الوسطى}؛ وتكون
الاستدلالات التي فوق~${\OLId{latin-looped}{S}}$ في~$\pi'$ كلها قضوية، والتي تحت~${\OLId{latin-looped}{S}}$ كلها
سوَرية. (ونسمي~!!^a{formula}~$!A$ \emph{سوَرية المقدمة}، أو نقول إنها
في \emph{الصيغة سوَرية المقدمة}، متى كانت على الصورة
\[\mathsf{Q}_{\OLMathNumeral{1}}{\OLId{latin-ordinary}{x}}_{\OLMathNumeral{1}}\dots\mathsf{Q}_{\OLId{latin-ordinary}{n}}{\OLId{latin-ordinary}{x}}_{\OLId{latin-ordinary}{n}}\,!B({\OLId{latin-ordinary}{x}}_{\OLMathNumeral{1}},\dots,{\OLId{latin-ordinary}{x}}_{\OLId{latin-ordinary}{n}})\]
حيث يكون كل~$\mathsf{Q}_{\OLId{latin-ordinary}{i}}$ إما~$\lforall$ وإما~$\lexists$.) وتؤدي
مبرهنة المتتالية الوسطى بدورها إلى مبرهنة هيربراند. وصورتها العامة عسرة
العبارة بعض العسر، وخلاصتها أن لكل~!!{provable}~!!{sentence}~$!A$ في
منطق الرتبة الأولى قضية تحصيل حاصل~$!A'$ يصح أن~!!{prove} منها~$!A$.
وأما صورتها الخاصة بـ!!{sentence}s من
الصورة~$\lexists[{\OLId{latin-ordinary}{x}}][!B({\OLId{latin-ordinary}{x}})]$، حيث~$!B$ خالية من الأسوار:

\begin{thm}[مبرهنة هيربراند]
لنفرض أن~$!B({\OLId{latin-ordinary}{x}})$ خالية من الأسوار وأن
$\Proves \lexists[{\OLId{latin-ordinary}{x}}][!B({\OLId{latin-ordinary}{x}})]$. عندئذ توجد حدود~${\OLId{latin-ordinary}{t}}_{\OLMathNumeral{1}}$, \dots,~${\OLId{latin-ordinary}{t}}_{\OLId{latin-ordinary}{n}}$
بحيث~$\Proves !B({\OLId{latin-ordinary}{t}}_{\OLMathNumeral{1}}) \lor \dots \lor !B({\OLId{latin-ordinary}{t}}_{\OLId{latin-ordinary}{n}})$.
\end{thm}

ونسمي~!!{formula}~$!B({\OLId{latin-ordinary}{t}}_{\OLMathNumeral{1}}) \lor \dots \lor !B({\OLId{latin-ordinary}{t}}_{\OLId{latin-ordinary}{n}})$
\emph{فصل هيربراند} لـ$\lexists[{\OLId{latin-ordinary}{x}}][!B({\OLId{latin-ordinary}{x}})]$. ولما كانت~$!B$ خالية من
الأسوار، كان الفصل كذلك. فإن كان~!!{provable}، فهو~!!{provable} ببرهان
خال من القطع، وله من ثم~!!a{proof} لا يستعمل إلا الاستدلالات القضوية؛
فثبت أنه تحصيل حاصل.

موضع العسر من مبرهنة هيربراند إثبات فصل هيربراند لكل
!!{provable}~!!{formula}~$\lexists[{\OLId{latin-ordinary}{x}}][!B({\OLId{latin-ordinary}{x}})]$. وأما العكس، وهو قابلية
$\lexists[{\OLId{latin-ordinary}{x}}][!B({\OLId{latin-ordinary}{x}})]$ للبرهنة متى كان لها فصل هيربراند، فهين؛ لأن فصل
هيربراند يستلزم~$\lexists[{\OLId{latin-ordinary}{x}}][!B({\OLId{latin-ordinary}{x}})]$. وهذا، على سبيل المثال،
!!a{proof} في~\Log{G3c} للحالة~${\OLId{latin-ordinary}{n}}={\OLMathNumeral{2}}$:
\[
\Axiom$!B({\OLId{latin-ordinary}{t}}_{\OLMathNumeral{1}}) \fCenter \lexists[{\OLId{latin-ordinary}{x}}][!B({\OLId{latin-ordinary}{x}})], !B({\OLId{latin-ordinary}{t}}_{\OLMathNumeral{1}}), !B({\OLId{latin-ordinary}{t}}_{\OLMathNumeral{2}})$
\Axiom$!B({\OLId{latin-ordinary}{t}}_{\OLMathNumeral{2}}) \fCenter \lexists[{\OLId{latin-ordinary}{x}}][!B({\OLId{latin-ordinary}{x}})], !B({\OLId{latin-ordinary}{t}}_{\OLMathNumeral{1}}), !B({\OLId{latin-ordinary}{t}}_{\OLMathNumeral{2}})$
\RightLabel{\LeftR{\lor}}
\BinaryInf$!B({\OLId{latin-ordinary}{t}}_{\OLMathNumeral{1}}) \lor !B({\OLId{latin-ordinary}{t}}_{\OLMathNumeral{2}}) \fCenter \lexists[{\OLId{latin-ordinary}{x}}][!B({\OLId{latin-ordinary}{x}})], !B({\OLId{latin-ordinary}{t}}_{\OLMathNumeral{1}}), !B({\OLId{latin-ordinary}{t}}_{\OLMathNumeral{2}})$
\RightLabel{\RightR{\lexists}}
\UnaryInf$!B({\OLId{latin-ordinary}{t}}_{\OLMathNumeral{1}}) \lor !B({\OLId{latin-ordinary}{t}}_{\OLMathNumeral{2}}) \fCenter \lexists[{\OLId{latin-ordinary}{x}}][!B({\OLId{latin-ordinary}{x}})], !B({\OLId{latin-ordinary}{t}}_{\OLMathNumeral{2}})$
\RightLabel{\RightR{\lexists}}
\UnaryInf$!B({\OLId{latin-ordinary}{t}}_{\OLMathNumeral{1}}) \lor !B({\OLId{latin-ordinary}{t}}_{\OLMathNumeral{2}}) \fCenter \lexists[{\OLId{latin-ordinary}{x}}][!B({\OLId{latin-ordinary}{x}})]$
\DisplayProof
\]

ولنستخرج فصل هيربراند من~!!a{proof}~$\pi$ لـ
$\Sequent \lexists[{\OLId{latin-ordinary}{x}}][!B({\OLId{latin-ordinary}{x}})]$. نبدأ بالحالة التي تجتمع فيها جميع
استدلالات~$\RightR{\lexists}$ في آخر~!!{proof} الخالي من القطع~$\pi$:
\[
\AxiomC{}
\RightLabel{$\pi_{\OLMathNumeral{1}}$}
\Deduce$ \fCenter \lexists[{\OLId{latin-ordinary}{x}}][!B({\OLId{latin-ordinary}{x}})], !B({\OLId{latin-ordinary}{t}}_{\OLMathNumeral{1}}), \dots, !B({\OLId{latin-ordinary}{t}}_{\OLId{latin-ordinary}{n}})$
\RightLabel{\RightR{\lexists}}
\UnaryInf$ \fCenter \lexists[{\OLId{latin-ordinary}{x}}][!B({\OLId{latin-ordinary}{x}})], !B({\OLId{latin-ordinary}{t}}_{\OLMathNumeral{2}}), \dots, !B({\OLId{latin-ordinary}{t}}_{\OLId{latin-ordinary}{n}})$
\RightLabel{$({\OLId{latin-ordinary}{n}}-{\OLMathNumeral{1}}) \times \RightR{\lexists}$}
\Deduce$ \fCenter \lexists[{\OLId{latin-ordinary}{x}}][!B({\OLId{latin-ordinary}{x}})]$
\DisplayProof
\]
فإن كانت هذه جميع استدلالات~$\RightR{\lexists}$ في~$\pi$ التي تكون
فيها~$\lexists[{\OLId{latin-ordinary}{x}}][!B({\OLId{latin-ordinary}{x}})]$ فعالة، لم يقع في~$\pi_{\OLMathNumeral{1}}$ استدلال سوَري آخر؛
وذلك أن~$!B({\OLId{latin-ordinary}{x}})$ لا سور فيها، وأن متتالية النهاية لا~!!{formula}s في
يسارها. وعندئذ تكون
$\fCenter \lexists[{\OLId{latin-ordinary}{x}}][!B({\OLId{latin-ordinary}{x}})], !B({\OLId{latin-ordinary}{t}}_{\OLMathNumeral{1}}), \dots, !B({\OLId{latin-ordinary}{t}}_{\OLId{latin-ordinary}{n}})$ المتتالية
الوسطى لـ$\pi$. وفوق ذلك، لما لم توجد استدلالات سوَرية فوق المتتالية
الوسطى، احتوتها كل المتتاليات الواقعة فوقها، ولم تكن فيها
$\lexists[{\OLId{latin-ordinary}{x}}][!B({\OLId{latin-ordinary}{x}})]$ فعالة ولا رئيسة. ومن ثم يمكن حذفها من~!!{proof}
$\pi_{\OLMathNumeral{1}}$، فنحصل على~!!a{proof} لـ
$\fCenter !B({\OLId{latin-ordinary}{t}}_{\OLMathNumeral{1}}), \dots, !B({\OLId{latin-ordinary}{t}}_{\OLId{latin-ordinary}{n}})$، وباستخدام~${\OLId{latin-ordinary}{n}}-{\OLMathNumeral{1}}$ تطبيقًا لـ
$\RightR{\lor}$ ينتج لنا~!!a{proof} لفصل هيربراند
$\Sequent !B({\OLId{latin-ordinary}{t}}_{\OLMathNumeral{1}}) \lor \dots \lor !B({\OLId{latin-ordinary}{t}}_{\OLId{latin-ordinary}{n}})$.

لكن لا يجوز لنا أن نفترض، حتى في~!!{proof} خال من القطع لـ
$\Sequent \lexists[{\OLId{latin-ordinary}{x}}][!B({\OLId{latin-ordinary}{x}})]$، أن استدلالات~$\RightR{\lexists}$ كلها
كتلة واحدة في آخره؛ فقد تتخللها استدلالات تكون فيها إحدى~$!B({\OLId{latin-ordinary}{t}}_{\OLId{latin-ordinary}{i}})$
رئيسة. ومن أجل هذا نحتاج إلى مبرهنة المتتالية الوسطى.

\begin{thm}[مبرهنة المتتالية الوسطى]\ollabel{thm:midsequent} إذا وُجد
  في~\Log{G3c}~!!{proof}~$\pi$ لـ${\OLId{greek-variable}{Gamma}} \Sequent {\OLId{greek-variable}{Delta}}$ تكون فيه كل
  !!{formula}s في~${\OLId{greek-variable}{Gamma}}$ و${\OLId{greek-variable}{Delta}}$ سوَرية المقدمة، وُجد~!!a{proof}
  $\pi'$ يحتوي متتالية~$\Pi \Sequent \Lambda$ بحيث تكون كل
  الاستدلالات تحتها سوَرية، وكل الاستدلالات فوقها قضوية.
\end{thm}

\begin{proof}
  عرّف \emph{رتبة} الاستدلال السُّوَري في~$\pi$ بأنها عدد الاستدلالات
  القضوية الواقعة تحته، وعرّف رتبة~$\pi$، أي~${\OLId{latin-ordinary}{o}}(\pi)$، بأنها مجموع رتب
  كل استدلالاته السُّوَرية. إذا كان~${\OLId{latin-ordinary}{o}}(\pi)={\OLMathNumeral{0}}$، لم يكن تحت أي استدلال
  سُوَري استدلال قضوي، فنكون قد انتهينا. فإن لم يقع في~$\pi$ استدلال
  سُوَري أصلًا، كانت المتتالية النهائية نفسها متتالية وسطى. وإلا، فلأن
  الاستدلالات السُّوَرية أحادية المقدمة، تؤلف الاستدلالات السُّوَرية
  الواقعة في أسفل البرهان سلسلة واحدة، ومقدمة أعلاها هي المتتالية الوسطى.

  وإذا كان~${\OLId{latin-ordinary}{o}}(\pi)>{\OLMathNumeral{0}}$، فاختر أدنى استدلال سُوَري رتبته~$>{\OLMathNumeral{0}}$. ولأن
  رتبته موجبة، فلا بد من وجود استدلال قضوي تحته. ولأنه أدنى استدلال
  سُوَري من هذا القبيل، لا يجوز أن تكون نتيجته مقدمة لاستدلال سُوَري؛
  وإلا كان ذلك الاستدلال السُّوَري الثاني الأدنى تحته استدلال قضوي
  كذلك. ونميز الحالات الكثيرة بحسب نوع الاستدلال السُّوَري ونوع
  الاستدلال القضوي الواقع تحته. ولأن متتالية النهاية لا تتألف إلا من
  صيغ سوَرية المقدمة، لا يجوز أن تكون الصيغة الرئيسة للاستدلال السُّوَري
  فعالة في الاستدلال القضوي التالي. وإلا كان لصيغة ذلك الاستدلال الرئيسة
  سور في نطاق رابط قضوي. وبخاصية~!!{formula} الجزئية، لزم أن تكون هذه
  !!{formula} جزئية من~!!a{formula} في متتالية النهاية. ومن ثم
  تكون~!!{formula} الرئيسة للاستدلال السُّوَري~!!a{element} من سياق
  الاستدلال القضوي الأدنى.

  وفيما يأتي مثالان:

أولًا، لنفرض أن لدينا~$\RightR{\lforall}$ ينتهي في المقدمة اليمنى لـ
$\LeftR{\lif}$. عندئذ ينتهي~!!{proof} الجزئي المعني من~$\pi$ بـ
!!{proof} الجزئي~$\pi_{\OLMathNumeral{1}}$:
\[
\AxiomC{}
\RightLabel{$\pi_{\OLMathNumeral{2}}$}
\Deduce${\OLId{greek-variable}{Gamma}}' \fCenter {\OLId{greek-variable}{Delta}}', \lforall[{\OLId{latin-ordinary}{x}}][!A({\OLId{latin-ordinary}{x}})], !B$
\AxiomC{}
\RightLabel{$\pi_{\OLMathNumeral{3}}$}
\Deduce$!C, {\OLId{greek-variable}{Gamma}}' \fCenter {\OLId{greek-variable}{Delta}}', !A({\OLId{latin-ordinary}{c}})$
\RightLabel{\RightR{\lforall}}
\UnaryInf$!C, {\OLId{greek-variable}{Gamma}}' \fCenter {\OLId{greek-variable}{Delta}}', \lforall[{\OLId{latin-ordinary}{x}}][!A({\OLId{latin-ordinary}{x}})]$
\RightLabel{\LeftR{\lif}}
\BinaryInf$!B \lif !C, {\OLId{greek-variable}{Gamma}}' \fCenter {\OLId{greek-variable}{Delta}}', \lforall[{\OLId{latin-ordinary}{x}}][!A({\OLId{latin-ordinary}{x}})]$
\DisplayProof
\]
يمكننا افتراض أن~$\pi$ منتظم، ومن ثم لا يظهر المتغير المميز~${\OLId{latin-ordinary}{c}}$ خارج
$\pi_{\OLMathNumeral{3}}$؛ وبوجه خاص لا يظهر في~$!B \lif !C$ ولا في~$\pi_{\OLMathNumeral{2}}$. والآن
تأمل البرهان~$\pi_{\OLMathNumeral{1}}'$:
\[
\AxiomC{}
\RightLabel{$\pi_{\OLMathNumeral{2}}'$}
\Deduce${\OLId{greek-variable}{Gamma}}' \fCenter {\OLId{greek-variable}{Delta}}', \lforall[{\OLId{latin-ordinary}{x}}][!A({\OLId{latin-ordinary}{x}})], !A({\OLId{latin-ordinary}{c}}), !B$
\AxiomC{}
\RightLabel{$\pi_{\OLMathNumeral{3}}$}
\Deduce$!C, {\OLId{greek-variable}{Gamma}}' \fCenter {\OLId{greek-variable}{Delta}}', \lforall[{\OLId{latin-ordinary}{x}}][!A({\OLId{latin-ordinary}{x}})], !A({\OLId{latin-ordinary}{c}})$
\RightLabel{\LeftR{\lif}}
\BinaryInf$!B \lif !C, {\OLId{greek-variable}{Gamma}}' \fCenter {\OLId{greek-variable}{Delta}}', \lforall[{\OLId{latin-ordinary}{x}}][!A({\OLId{latin-ordinary}{x}})], !A({\OLId{latin-ordinary}{c}})$
\RightLabel{\RightR{\lforall}}
\UnaryInf$!B \lif !C, {\OLId{greek-variable}{Gamma}}' \fCenter {\OLId{greek-variable}{Delta}}', \lforall[{\OLId{latin-ordinary}{x}}][!A({\OLId{latin-ordinary}{x}})], \lforall[{\OLId{latin-ordinary}{x}}][!A({\OLId{latin-ordinary}{x}})]$
\DisplayProof
\]
حيث ينتج~$\pi_{\OLMathNumeral{2}}'$ من~$\pi_{\OLMathNumeral{2}}$ بإضعافه بـ$!A({\OLId{latin-ordinary}{c}})$ في اليمين. (ولا يضيف
ذلك أي استدلال، وبوجه خاص لا يضيف استدلالات سوَرية قد تؤثر في الرتبة.
كما أنه لا يخل بأي شرط لمتغير مميز في~$\pi_{\OLMathNumeral{2}}$، إذ لا يُستخدم أي ثابت
في~$!A({\OLId{latin-ordinary}{c}})$ متغيرًا مميزًا في~$\pi_{\OLMathNumeral{2}}$.) ويظل شرط المتغير المميز في
$\RightR{\lforall}$ متحققًا، لأن~${\OLId{latin-ordinary}{c}}$ لا تظهر في~$!B \lif !C$.

نستدعي مقبولية التقليص فنحصل على~!!a{proof}~$\pi_{\OLMathNumeral{1}}''$ لـ
$!B \lif !C, {\OLId{greek-variable}{Gamma}}' \fCenter {\OLId{greek-variable}{Delta}}', \lforall[{\OLId{latin-ordinary}{x}}][!A({\OLId{latin-ordinary}{x}})]$. ولم
يغير برهان مقبولية~$\RightR{\Contraction}$ لـ$\lforall[{\OLId{latin-ordinary}{x}}][!A({\OLId{latin-ordinary}{x}})]$،
ولا برهان قابلية~$\RightR{\lforall}$ للعكس المستخدم فيه، رتبة أي
استدلال؛ فقد حذفا استدلالات في أقصى الأحوال. ولذلك لا توجد استدلالات
سوَرية في~$\pi_{\OLMathNumeral{1}}''$ أكثر مما في~$\pi_{\OLMathNumeral{1}}$، ولكل استدلال سُوَري في
$\pi_{\OLMathNumeral{1}}''$ عدد الاستدلالات القضوية نفسه تحته كما للاستدلال المقابل في
$\pi_{\OLMathNumeral{1}}$، باستثناء~$\RightR{\lforall}$ الأخير: كان تحته
$\LeftR{\lif}$ في~$\pi$، وقد نقلناه إلى فوقه في~$\pi_{\OLMathNumeral{1}}''$. استبدل
$\pi_{\OLMathNumeral{1}}''$ بـ$\pi_{\OLMathNumeral{1}}$ في~$\pi$. ويكون للـ!!{proof} الناتج رتبة أدنى،
لأن رتبة استدلال~$\RightR{\lforall}$ على الأقل نقصت بمقدار~${\OLMathNumeral{1}}$. ثم
طبّق فرض الاستقراء.

والحالات التي يكون فيها الاستدلال السُّوَري هو~$\RightR{\lexists}$
أسهل، إذ لا نحتاج إلى التقليص ولا إلى القلق بشأن شرط متغير مميز. فافترض
مثلًا أن~$\RightR{\lexists}$ يتبعه~$\RightR{\land}$، وهذه المرة بوصفه
المقدمة اليسرى. يكون~!!{proof} الجزئي المعني هو~$\pi_{\OLMathNumeral{1}}$:
\[
\AxiomC{}
\RightLabel{$\pi_{\OLMathNumeral{2}}$}
\Deduce${\OLId{greek-variable}{Gamma}}' \fCenter {\OLId{greek-variable}{Delta}}', \lexists[{\OLId{latin-ordinary}{x}}][!A({\OLId{latin-ordinary}{x}})], !A({\OLId{latin-ordinary}{t}}), !B$
\RightLabel{\RightR{\lexists}}
\UnaryInf${\OLId{greek-variable}{Gamma}}' \fCenter {\OLId{greek-variable}{Delta}}', \lexists[{\OLId{latin-ordinary}{x}}][!A({\OLId{latin-ordinary}{x}})], !B$
\AxiomC{}
\RightLabel{$\pi_{\OLMathNumeral{3}}$}
\Deduce${\OLId{greek-variable}{Gamma}}' \fCenter {\OLId{greek-variable}{Delta}}', \lexists[{\OLId{latin-ordinary}{x}}][!A({\OLId{latin-ordinary}{x}})], !C$
\RightLabel{\RightR{\land}}
\BinaryInf${\OLId{greek-variable}{Gamma}}' \fCenter {\OLId{greek-variable}{Delta}}', \lexists[{\OLId{latin-ordinary}{x}}][!A({\OLId{latin-ordinary}{x}})], !B \land !C$
\DisplayProof
\]
استبدل~$\pi_{\OLMathNumeral{1}}'$ بـ$\pi_{\OLMathNumeral{1}}$ في~$\pi$:
\[
\AxiomC{}
\RightLabel{$\pi_{\OLMathNumeral{2}}$}
\Deduce${\OLId{greek-variable}{Gamma}}' \fCenter {\OLId{greek-variable}{Delta}}', \lexists[{\OLId{latin-ordinary}{x}}][!A({\OLId{latin-ordinary}{x}})], !A({\OLId{latin-ordinary}{t}}), !B$
\AxiomC{}
\RightLabel{$\pi_{\OLMathNumeral{3}}'$}
\Deduce${\OLId{greek-variable}{Gamma}}' \fCenter {\OLId{greek-variable}{Delta}}', \lexists[{\OLId{latin-ordinary}{x}}][!A({\OLId{latin-ordinary}{x}})], !A({\OLId{latin-ordinary}{t}}), !C$
\RightLabel{\RightR{\land}}
\BinaryInf${\OLId{greek-variable}{Gamma}}' \fCenter {\OLId{greek-variable}{Delta}}', \lexists[{\OLId{latin-ordinary}{x}}][!A({\OLId{latin-ordinary}{x}})], !A({\OLId{latin-ordinary}{t}}), !B \land !C$
\RightLabel{\RightR{\lexists}}
\UnaryInf${\OLId{greek-variable}{Gamma}}' \fCenter {\OLId{greek-variable}{Delta}}', \lexists[{\OLId{latin-ordinary}{x}}][!A({\OLId{latin-ordinary}{x}})], !B \land !C$
\DisplayProof
\]
حيث ينتج~$\pi_{\OLMathNumeral{3}}'$ من~$\pi_{\OLMathNumeral{3}}$ بإضعافه بـ$!A({\OLId{latin-ordinary}{t}})$ في اليمين. ولا يتضمن
هذا الإضعاف إلا إضافة~$!A({\OLId{latin-ordinary}{t}})$ إلى يمين كل متتالية في~$\pi_{\OLMathNumeral{3}}$، فلا يغير
رتبة أي استدلال سُوَري. وتكون رتب الاستدلالات السُّوَرية في~$\pi'$
كما هي في~$\pi$، سوى أن رتبة استدلال~$\RightR{\lexists}$ الذي بُدّل
موضعه مع~$\RightR{\land}$ نقصت بمقدار~${\OLMathNumeral{1}}$. فينطبق فرض الاستقراء.
\end{proof}

\begin{prob}
صف الحالتين في برهان~\olref{thm:midsequent} اللتين ينتهي في إحداهما
$\LeftR{\lexists}$ في مقدمة~$\RightR{\lif}$، وينتهي في الأخرى
$\LeftR{\lforall}$ في المقدمة اليسرى لـ$\LeftR{\lor}$.
\end{prob}

\begin{proof}[برهان مبرهنة هيربراند]
لنفرض أن~$\pi$ ينتهي بـ$\Sequent \lexists[{\OLId{latin-ordinary}{x}}][!B({\OLId{latin-ordinary}{x}})]$. وبفضل
\olref{thm:midsequent} نستطيع تحويل~$\pi$ إلى~!!a{proof}~$\pi'$ ذي
متتالية وسطى~${\OLId{latin-looped}{S}}$، ولا بد أن تكون على الصورة
\[
\Sequent \lexists[{\OLId{latin-ordinary}{x}}][!B({\OLId{latin-ordinary}{x}})], !B({\OLId{latin-ordinary}{t}}_{\OLMathNumeral{1}}), \dots, !B({\OLId{latin-ordinary}{t}}_{\OLId{latin-ordinary}{n}}).
\]
ولأن كل الاستدلالات فوق~${\OLId{latin-looped}{S}}$ قضوية، لا يجوز أن تكون
$\lexists[{\OLId{latin-ordinary}{x}}][!B({\OLId{latin-ordinary}{x}})]$ رئيسة في استدلال فوق~${\OLId{latin-looped}{S}}$. ولأنها ليست ذرية، لا
يجوز أن تكون رئيسة في مسلّمة. ولأن~$\lexists[{\OLId{latin-ordinary}{x}}][!B({\OLId{latin-ordinary}{x}})]$ لا يجوز أن تكون
صيغة جزئية حقيقية من~$!B({\OLId{latin-ordinary}{t}}_{\OLMathNumeral{1}})$، إذ إن~$!B({\OLId{latin-ordinary}{x}})$ خالية من الأسوار، فلا
يمكن كذلك أن تكون فعالة فوق~${\OLId{latin-looped}{S}}$. ومن ثم يعطي حذف
$\lexists[{\OLId{latin-ordinary}{x}}][!B({\OLId{latin-ordinary}{x}})]$ من~!!{proof} الجزئي المنتهي بـ${\OLId{latin-looped}{S}}$~!!{proof}
صحيحًا لـ$\Sequent !B({\OLId{latin-ordinary}{t}}_{\OLMathNumeral{1}}), \dots, !B({\OLId{latin-ordinary}{t}}_{\OLId{latin-ordinary}{n}})$، ومنه يمكن الحصول على
!!a{proof} لـ$\Sequent !B({\OLId{latin-ordinary}{t}}_{\OLMathNumeral{1}}) \lor \dots \lor !B({\OLId{latin-ordinary}{t}}_{\OLId{latin-ordinary}{n}})$ بتطبيق
$\RightR{\lor}$ عدد~${\OLId{latin-ordinary}{n}}-{\OLMathNumeral{1}}$ من المرات.
\end{proof}

% \begin{prob}
% أوجد فصل هيربراند لـ$\lexists[x][\lforall[y][(!A(x) \lif
% !A(y))]]$ بإيجاد~!!{proof} خال من القطع لـ$\Sequent
% \lexists[x][\lforall[y][(!A(x) \lif !A(y))]]$ وتطبيق التحويل في
% \olref{thm:midsequent} لإيجاد متتالية وسطى عند الحاجة.
% \end{prob}

\end{document}
